\subsection{A 104-donor deduplicated single-cell atlas contains 117 prespecified gene programs}\label{a-102-donor-single-cell-atlas-contains-117-prespecified-gene-programs}

To place disease-associated gene programs in specific liver cell types, we integrated seven single-cell datasets into one annotated atlas (Ramachandran et al., 2019; Andrews et al., 2022; Filliol et al., 2022; Guilliams et al., 2022; Gribben et al., 2024; Kim et al., 2025) (Figure 4A). After deduplicating repeated sequencing captures and removing 41 antibody-tag rows, the unique-cell atlas comprises 650,083 unique cells across 145 physical captures and 104 source donors (with 226,457 unsorted unique cells across 63 donors entering unsorted donor-level analyses; historical run-level inputs used to discover the frozen weights held 1,232,318 cells across 275 libraries). The 117 prespecified gene programs (DeTomaso and Yosef, 2021) discovered in five lineage subsets remain frozen in membership and weight. We summarized program, communication and regulator analyses at the donor level, because sequencing runs, sort fractions and cells are not biological replicates. Source-qualified descriptive categories are 17 Cirrhosis, 23 Healthy, 36 Steatohepatitis, 14 Steatosis, 12 Unresolved. Deposited diagnosis, measured histology and etiology were kept separate. GSE202379 P98 was excluded from disease inference because its Healthy-control label conflicts with SAF S1A1F1. P30 (S0A1F0) was retained as a qualified nonsteatotic reference, with an exclusion sensitivity; it is not confirmed histologically normal tissue. Official RNA/GSM metadata place local GSE244832 D11/JB\_\allowbreak{}321 in MASL and D14/JB\_\allowbreak{}337 in MASH. This correction does not verify RNA--ATAC pair identities. All 19 GSE192742 donors remain outside disease inference. Alcohol/PBC cirrhosis is not classified as MASLD, and GSE189600 replicate independence remains unresolved.

The atlas follows the same evidence rules as the rest of the Resource. A program received an evidence state only when donor identity, retained gene weight, cell-state coverage and complete multiple-testing families were all in place. If any one failed, the program was untestable. Untestable is not evidence that the biological program is absent. Each selected program keeps its unresolved alternative and its next discriminating experiment. The integrated UMAP shows 50,000 cells sampled uniformly from the analyzed cells and colored by cell type (Figure 4B). It is descriptive, and cells are not the unit of inference. We do not draw a donor-averaged disease density, because its Healthy and Steatohepatitis groups come from different mixes of datasets.

\begin{mainfigureplate}{4}{468bp}
\noindent\makebox[\linewidth][c]{\paperpanel{A}{132bp}{figures/main/figure4/fig4a_atlas_program_schematic.pdf}\hspace{8bp}\paperpanel{B}{303bp}{figures/main/figure4/fig4b_scrna_umap_embeddable.pdf}}
\par\vspace{8pt}
\noindent\makebox[\linewidth][c]{\paperpanel{C}{180bp}{figures/main/figure4/fig4c_same_atlas_lineage_specificity.pdf}\hspace{8bp}\paperpanel{D}{170bp}{figures/main/figure4/fig4d_steatosite_program_composition.pdf}}
\par\vspace{8pt}
\noindent\makebox[\linewidth][c]{\paperpanel{E}{265bp}{figures/main/figure4/fig4e_hepatocyte_trio_forest.pdf}\hspace{8bp}\paperpanel{F}{140bp}{figures/main/figure4/fig4f_ccl20_state_set_excess.pdf}}
\par\vspace{8pt}
\noindent\makebox[\linewidth][c]{\paperpanel{G}{255bp}{figures/main/figure4/fig4g_state_gene_hepatocyte_share.pdf}\hspace{8bp}\paperpanel{H}{180bp}{figures/main/figure4/fig4h_pnpla3_inflamed_score.pdf}}

\end{mainfigureplate}
\papercaption{4}


\subsection{Gene content contradicts the discovery lineage for 39 of 117 programs}\label{gene-content-contradicts-the-discovery-lineage-for-37-of-117-programs}

A program's name records the cell-type subset in which it was discovered, not what its genes are. To test whether the genes match the name, we scored each program against deduplicated donor-level expression across the six active unsorted lineages (excluding sorted GSE136103 cells). Content matched the declared lineage for 78 of 117 programs (66.7\%, 95\% donor-bootstrap CI 65.0\% to 67.5\%) and contradicted it for 39 (Figure 4C). The mismatches are systematic: 18 programs discovered outside hepatocytes (4 fibroblast, 9 cholangiocyte, 4 macrophage, 1 T-cell) have hepatocyte content, and several are named for it. Examples are a ``Hepatocyte xenobiotic (CYP2C)'' program labelled fibroblast and a ``Bile-acid transport (ABCB11)'' program labelled cholangiocyte. These hepatocyte-content mismatches carry a higher measured burden of ambient RNA, the free transcripts captured along with each cell, than correctly labelled non-hepatocyte programs (median 0.432 against 0.061). Hepatocyte transcripts are the dominant ambient species in these data. This association does not show that ambient RNA explains every mismatch. The panels separate three lineages for each program: the lineage its content matches, the lineage it was discovered in and the lineage used for disease testing. The program definitions are unchanged.

In separate donor-paired contrasts within the same atlas, Stromal ECM/IGFBP7 scored highest in fibroblasts (beta = 2.271 over hepatocytes; 53 donors) and Ductular injury/BICC1 scored highest in cholangiocytes (beta = 1.574 over hepatocytes; 51 donors; Figure S4A,B). Both pass the ten-contrast BH and HC3 families (Long and Ervin, 2000). This is descriptive localization. It does not show where the transcripts came from or that the program acts within that cell type.

\subsection{Most programs change with fibrosis stage in 493 held-out livers, mainly with cell mix}\label{most-programs-change-with-fibrosis-stage-in-493-held-out-livers}

The discovery cohorts could not resolve stage associations of the programs. At most 54 unique donors entered the single-cell test, and no program passed correction (smallest BH q = 0.232; Figure S4AB). Three bulk cohorts gave no fibrosis or NAS call over 224 tests, with t on 2 degrees of freedom (Figure S4AC). These analyses, the continuum profiles and the 43 molecular systems are in Supplementary Note 20 and Figures S4AB to S4AF. We therefore scored the 117 fixed programs in SteatoSITE, which no part of the program map used (493 staged donors; Supplementary Note 17). Of 112 testable programs, 98 changed with fibrosis stage (BH q \textless{} 0.05 over 112; 63 higher and 35 lower with stage; Figure 4D). Both prespecified confirmatory programs passed unadjusted (Holm over two): Stromal ECM rose by 0.674 and Ductular injury by 0.633 score SD per stage step. Adding seven isometric log-ratio cell-composition coordinates from 22-autosome leave-one-chromosome-out (LOCO) deconvolution separating hepatocytes and cholangiocytes moved most programs toward zero (only 4 of 112 programs retained BH q \textless{} 0.05 across all 22 LOCO chromosomes; Figure 4D). Stromal ECM survived composition adjustment across all 22 LOCO chromosomes (modal chr3 $\beta$ = 0.0881, BH q = 8.93 $\times$ 10$^{-7}$), whereas Ductular injury attenuated under modal-chr2 LOCO ($\beta$ = 0.0297, BH q = 0.402; 0 of 22 chromosomes at q \textless{} 0.05; median $\beta$ = 0.0454). In the discovery bulk cohorts, fibrosis estimates were positive for 20 of 21 cholangiocyte-content and 19 of 21 fibroblast-content programs and negative for 29 of 36 hepatocyte-content programs (descriptive). Bulk program scores therefore track stage mainly together with more bile-duct and fibroblast RNA and less hepatocyte RNA. Even with LOCO exclusion, the composition terms are deconvolution estimates from the same libraries, so the adjusted coefficients are descriptive and do not measure change inside cells. The next three analyses look inside hepatocytes directly.

\subsection{GNMT, MAT1A and CYP2C19 are lower per hepatocyte}\label{gnmt-mat1a-and-cyp2c19-are-lower-per-hepatocyte}

The three methionine-cycle and drug-metabolism genes that are lower in bulk RNA and protein (Figure S3X) were tested in hepatocyte pseudobulk from five single-nucleus cohorts with one HC3-studentized model (Figure 4E). In GSE289173, a cohort not used in any earlier analysis of these genes (23 control and 17 steatohepatitis donors after excluding GSM8785704), all three were lower per hepatocyte: CYP2C19 -1.69 log2 CPM (95\% CI -2.56 to -0.81; Holm \emph{p} = 0.0034), GNMT -1.17 (-1.92 to -0.41; Holm \emph{p} = 0.0045) and MAT1A -0.92 (-1.44 to -0.41; Holm \emph{p} = 0.0034). Control and disease livers in that cohort were collected as separate series. GSE202379 and GSE244832, with 8 healthy and 44 disease deduplicated donors, were uninformative (GNMT Holm \emph{p} = 0.750, MAT1A \emph{p} = 0.691, CYP2C19 \emph{p} = 0.691). The full analysis is described with the bulk results.

\subsection{A CCL20+ hepatocyte state with an \texorpdfstring{NF-$\kappa$B}{NF-kB} response recurs across cohorts}\label{a-ccl20-hepatocyte-state-with-an-nf-kb-response-recurs-across-cohorts}

CCL20, TRIM31 and FCAMR are among the largest increases in the 774-unit bulk fit (12th, 30th and 48th of 6,390 genes higher in disease); their bulk rise with NAS was reported before (Hoang et al., 2019). In deduplicated MASLD donors, cells labeled hepatocyte held medians of 0.75 (GSE298719), 0.88 (GSE244832), 0.90 (GSE256398) and 0.50 (GSE202379) of CCL20 UMIs, and 0.97 to 1.00 of FCAMR and TRIM31 UMIs (Figure 4G). In single-nucleus data some of this share is ambient RNA, so these values are upper bounds. We then asked whether the CCL20+ cells labeled hepatocyte are hepatocytes by expression (unit: deduplicated donor; prespecified). After removing cells carrying PTPRC or non-hepatocyte markers, hepatocyte-assigned cells still held medians of 0.61 (GSE244832), 0.61 (GSE256398), 0.34 (GSE298719) and 0.33 (GSE202379) of all CCL20 UMIs; in the strict contamination-cleaned identity test, GSE244832 passed (7 donors, Holm \emph{p} = 0.0078), GSE256398 was non-significant (6 donors, Holm \emph{p} = 0.219), and GSE298719 and GSE202379 had fewer than 5 eligible donors after absolute filtering. Most CCL20 RNA in the larger snRNA livers therefore sits in cells with hepatocyte expression profiles, though no cohort has doublet scores and ductular cells with low KRT19 and KRT7 would pass as hepatocytes. In GSE298719 and GSE256398, CCL20+ hepatocytes expressed NF-$\kappa$B target genes such as CXCL1, CXCL8, BIRC3, ICAM1, TNFAIP3 and RELB more than CCL20$-$ hepatocytes of the same donor matched for sequencing depth and an injury score. We fixed this 64-gene set (62 measured) and tested it in two deduplicated cohorts not used to define it. The set was higher in CCL20+ hepatocytes in all 8 eligible GSE244832 donors (mean +0.85 log2; exact sign-flip \emph{p} = 0.0078, Holm \emph{p} = 0.0156 over two cohorts) and in all 5 GSE202379 donors (mean +1.01; \emph{p} = 0.0625, the smallest value possible with five donors) (Figure 4F). It could not be tested in GSE289173, where only 1 donor passed the 10-cell CCL20+ threshold. These analyses are exploratory: the replication test was written after the discovery cohorts were read, although the gene set was fixed before the two test cohorts were. Hepatocytes with NF-$\kappa$B and RELB activity associated with fibrosis have been described in human MASLD (Elison et al., 2025), as have senescence-associated and disease-associated hepatocyte states (Du et al., 2025; Qiu et al., 2026), so the state is not new; what these data add is that a fixed version of this response recurs in donors from other cohorts, and that most CCL20 RNA in human MASLD tissue is in cells with hepatocyte expression profiles. Our search for earlier human assignments of the CCL20 source was not systematic.

In five bulk cohorts, a five-gene score of this state (FCAMR, TRIM31, CXCL10, IL32, CCL20) rose by 0.153 SD per PNPLA3 148M allele after adjusting for cohort, inferred sex and four RNA-called genetic ancestry PCs (95\% CI 0.070 to 0.236; \emph{p} = 3.26 $\times$ 10$^{-4}$; n = 757; European-ancestry subset with 4 PCs +0.147 SD, 95\% CI 0.059 to 0.236, n = 641; unadjusted without PCs +0.176 SD, \emph{p} = 2.0 $\times$ 10$^{-5}$; Figure 4H). In the 633 participants with clean histology (excluding unstated GSE213621 Control fibrosis codes), adding histology and 4 ancestry PCs removed the effect (+0.029 SD, 95\% CI -0.045 to 0.102, \emph{p} = 0.441; EUR + clean histology + 4 PCs +0.026 SD, n = 547, \emph{p} = 0.510). The allele raises the score about as much as its effect on histology predicts, so these data do not show a genotype effect on the state beyond disease severity. Genotypes were called from RNA reads, and GSE135251 rs738409 shows modest departure from Hardy--Weinberg equilibrium (\emph{p} = 0.018 overall, 0.032 in EUR).

\subsection{Program scores covary across bulk livers in two cohorts not used to build the program map}\label{program-scores-covary-across-bulk-livers-in-two-cohorts-not-used-to-build-the-program-map}

We next asked whether program scores covary with one another across bulk livers (Supplementary Note 10). In the four discovery cohorts, 2,426 of 6,195 tested pairs of the 113 bulk-scorable programs had a positive partial correlation that passed BH correction, after adjustment for sex, fibrosis stage and NAS and after removal of the leading residual component. Of these, 650 pairs join programs of the same lineage and 1,776 join programs of different lineages. This discovery step used 469 source records from GSE130970, GSE135251, GSE162694 and GSE174478, scored on counts that precede the identity-qualified 774-unit dataset. Its unit is therefore not the genetic-identity cluster used in the other bulk analyses. In bulk tissue, covariance of program scores mixes cell-type abundance, biopsy content and program activity, so it does not show that cell types coordinate. A prespecified test in 258 GSE213621 disease units (genetic-identity clusters) used 2,328 pairs rediscovered on 371 identity-qualified discovery units, and it was untestable. Its program-group permutation control keeps covariance within each program and removes it between programs. This control still gave a mean excess over the non-coordinated baseline of 0.025 (95th percentile 0.081; prespecified limits 0.02 and 0.05), so the run stopped before any GSE213621 outcome statistic was computed. An earlier GSE193066 comparison rests on participant identities that cannot be qualified and is not counted as replication (Supplementary Note 10).

We then tested the 2,426 fixed discovery pairs, with a separate prespecified test and null, in two bulk cohorts that were not used to build or test the program map: GSE281797 (91 participants) and GSE225740 (87 participants), one biopsy each (Zhao et al., 2026; Powell et al., 2023). In both cohorts, the fixed pairs had positive histology-adjusted partial correlations more often than pairs placed at random among programs of the same lineages. This held after removing 1 and after removing 3 residual components. In each cohort the largest of the permutation and bootstrap \emph{p} values was 0.000999, the floor set by 1,000 permutations and 1,000 bootstrap draws, and the null was rejected in both cohorts under Holm correction over the two. With 1 component removed, the fraction of pairs with a positive correlation exceeded the random-placement mean by 0.158 in GSE281797 (0.848 against 0.690; participant-bootstrap 95\% interval of the excess 0.132 to 0.182) and by 0.312 in GSE225740 (0.834 against 0.522; 0.270 to 0.323) (Supplementary Note 10). The program-group permutation control that stopped the GSE213621 test was not run in these two cohorts. The two excesses are about 6 and 12 times the mean excess that control gave in GSE213621 (0.025), but that control used a different statistic and baseline, so the comparison is approximate. Participant overlap with the discovery cohorts was not checked by genotype in either cohort. GSE225740 identities cannot be checked at all, because it has no reads and no deposited sex, age or batch. GSE281797 is therefore the stronger replication, because it has full covariates and verifiable sex. In GSE225740, the result held after we added three estimated composition components, log library size and the top-100-gene count share to the design. The GSE281797 composition estimate failed its prespecified quality rule and was not read, so the GSE281797 result has no composition adjustment. None of these numbers is pooled with any GSE213621 or GSE193066 value. Estimated composition is not measured composition, so these results do not show that the covariance is independent of cell-type abundance.

To test whether the same pairs covary within tissue, we used Visium spots in 32 HMSMA livers (Li et al., 2025) and the 44 programs measurable in every array. Each donor contributed one array. Three arrays inferred to come from controls were excluded. The source study lists 7 controls among its 35 individuals without a key to the arrays, so 4 of the 32 analyzed donors are probably controls that cannot be identified. Program pairs with positive bulk covariance co-vary across spots, but coordinated pairs do not exceed other positively covarying pairs (134 against 80 pairs; excess 0.0116 Fisher z, \emph{p} = 0.124; an excess of 0.0280 or more would have been detected with 80\% power). Among cross-lineage pairs that share no genes, the 160 coordinated pairs correlated more across spots than the 372 non-coordinated pairs of the same lineage pairs (excess 0.0257 Fisher z; Holm \emph{p} = 0.0002; positive in 31 of 32 livers) (Supplementary Note 10). However, 292 of those 372 non-coordinated pairs have a gene-disjoint bulk correlation of zero or below. This difference therefore tracks the sign of bulk covariance, not coordination beyond the leading residual component. Each 55-micron spot holds several cells, so this is co-location at spot scale. It does not show expression in different cells or contact between cells. HMSMA has no key linking arrays to clinical records, so it carries no stage information. Neither the bulk nor the spot results show coordination between cell types.

Taken together, these analyses place cross-sectional stage-associated remodeling in donor-aware cell-type context and separate cell mix from change inside hepatocytes where the data allow. They do not redefine the program registry, and they do not treat histologic groups as states that change over time within a person. Colocalization gives no positive genetic call for IGFBP7 or BICC1, and the unresolved ambient-RNA sensitivity prevents assigning either program to hepatocytes from these data alone. We next asked whether the disease-state RNA changes appear in liver protein and in which cell types of intact tissue the genes are made (Figure 5).
